% Generated draft: requires mathematical and source review.
\documentclass[11pt]{article}
\usepackage[margin=25mm]{geometry}
\usepackage{fontspec}
\setmainfont{Latin Modern Roman}
\usepackage{amsmath,amssymb,mathtools,unicode-math}
\setmathfont{Latin Modern Math}
\usepackage{enumitem,needspace,xurl,hyperref}
\setlength{\parindent}{0pt}
\setlength{\parskip}{5pt}
\setlength{\emergencystretch}{3em}
\newcommand{\N}{\mathbb{N}}
\newcommand{\Z}{\mathbb{Z}}
\newcommand{\Q}{\mathbb{Q}}
\newcommand{\R}{\mathbb{R}}
\newcommand{\C}{\mathbb{C}}
\newcommand{\F}{\mathbb{F}}
\newcommand{\A}{\mathbb{A}}
\newcommand{\PP}{\mathbb{P}}
\newcommand{\E}{\mathbb{E}}
\newcommand{\eps}{\varepsilon}
\newcommand{\dd}{\,\mathrm{d}}
\begin{document}
\section*{UnsolvedMath ID: 7000007}
\textbf{Source problem code:} catalogue entry
\subsection{Definitions and mathematical statement}
In the usual Euclidean-surface setting, let \(i:S\hookrightarrow\R^3\) be a connected surface with induced metric \(g=i^*g_{\mathrm{Euc}}\), and let \(F\subset S\) be finite. A global rigidity interpretation asks whether every admissibly regular isometric immersion
\[
f:(S\setminus F,g)\to\R^3,\qquad f^*g_{\mathrm{Euc}}=g,
\]
is of the form \(f=Q\circ i+b\) on \(S\setminus F\), with \(Q\in O(3)\), \(b\in\R^3\). If only embeddings or only continuous isometric deformations are admitted, that is a different formulation. Infinitesimal rigidity, concerning linearized isometry fields, is another distinct question and is not silently substituted.

Ghomi asks whether some nonconvex surfaces retain the intended rigidity after deletion of finitely many points, particularly analytic tight surfaces and a torus of revolution. For embedded compact surfaces, a customary tightness convention is the two-piece property: intersection with every open half-space is connected or empty. A concrete analytic ring-torus model is
\[
i(u,v)=((R+r\cos v)\cos u,(R+r\cos v)\sin u,r\sin v),
\qquad R>r>0,
\]
with \(u,v\) periodic modulo \(2\pi\). The source does not say that every surface of revolution is intended or fix a particular generating curve.

The strongest finite-puncture existence reading is
\[
\exists S\text{ nonconvex}\ \forall F\subset S\text{ finite},
\qquad S\setminus F\text{ remains rigid}.
\]
Existence of one nonempty finite puncture set, or rigidity after every single puncture, gives weaker assertions; the source does not explicitly choose the quantifiers. It also leaves surface compactness, ambient category, competitor regularity/analyticity, and rigidity notion unspecified. The analytic regularity of the original surface does not by itself impose analyticity on all competitors.

The supplied research summary describes the punctured tight/nonconvex question as unresolved and mentions convex rigidity and possible puncture-flexibility phenomena without exact hypotheses. Those incomplete comparisons are not adopted as universal convex-surface theorems or as a resolution of the torus case.
\subsection{Short English statement}
Can a nonconvex surface remain isometrically rigid after finite punctures, especially an analytic tight surface such as a ring torus? The rigidity category and the quantifier over deleted points are essential missing conventions.
\subsection{Sources}
\begin{itemize}
\item Local UnsolvedMath registry, permanent ID \(7000007\), exact source code \texttt{catalogue entry}, ``Geometry of Curves and Surfaces --- Problem 1.7.'' Proposer Mohammad Ghomi, 2019.
\item Ghomi, \emph{Open Problems in Geometry of Curves and Surfaces} (2019), Problem 1.7, recorded PDF page 8; \url{https://people.math.gatech.edu/~ghomi/Papers/op.pdf}. Recorded access 2026-07-29; pdf-text-normalized extraction.
\item Supplied OPEN-TRIAGE research summary mentions Cohn--Vossen-type convex rigidity and puncture/nonrigidity literature, without complete references or matching regularity hypotheses. Its 2019--2026 no-resolution statement is a supplied search assessment.
\end{itemize}
\end{document}
